% PEARL Thesis -- Abstract
\chapter*{Abstract}
\addcontentsline{toc}{chapter}{Abstract}

Polycystic Ovary Syndrome (PCOS) is a heterogeneous endocrine disorder whose
ultrasound manifestations are difficult to assess consistently because images
contain speckle noise, low contrast, and acquisition-dependent variation. This
thesis presents \PEARL{}, a trust-oriented deep-learning pipeline for binary
PCOS classification from ovarian ultrasound images. The work focuses not only
on discrimination, but also on cross-dataset generalisation, calibration,
uncertainty estimation, and visual explanation.

The executed pipeline first evaluated nine ImageNet-pretrained backbones under
SRAD and Gaussian denoising on the deduplicated Figshare PCOS ultrasound
dataset. The strongest foundation models were then adapted to the PCOSgen
cohort using a train-pool/held-out-test protocol. The final pre-HPO ensemble
averages probabilities from DenseNet-121, ConvNeXt-Tiny, and ViT-B/16. On the
held-out PCOSgen test set of 1,468 images, the ensemble achieved an AUC of
0.9487 (95\% CI: 0.9368--0.9596), F1 of 0.8665, MCC of 0.7884, and Brier score
of 0.0981. Two-pass temperature scaling reduced ensemble ECE from 0.0854 to
0.0810. MC-Dropout with 50 stochastic passes provided predictive-entropy
summaries, while Grad-CAM exposed representative correct, uncertain, and
failure cases.

The results also expose a substantial domain-shift effect. Foundation
checkpoints performed strongly on the internal Figshare split but collapsed on
zero-shot PCOSgen evaluation. Fine-tuning on the PCOSgen training pool
recovered external AUC into the 0.93 range, and probability averaging improved
the final result further. The thesis documents the complete methodology,
including deduplication, split construction, SRAD, augmentation, model
training, HPO, calibration, uncertainty, explainability, reproducibility, and
the difference between features designed in the pipeline specification and
features actually used in the reported run.

\textbf{Keywords:} Polycystic Ovary Syndrome, ovarian ultrasound, transfer
learning, cross-dataset generalisation, Speckle-Reducing Anisotropic Diffusion,
calibration, MC-Dropout, Grad-CAM, uncertainty estimation, ensemble learning.
